%% kapa: exam=1 round=36 year=2025 period=1 form=A subject=civil_law subject_ko=민법 seq=37 number=37
%% source: Q-Net 감정평가사 기출문제 (cst00309) artlSeq=5231525 file=raw/1008/5231525/2025년 제36회 감정평가사 1차시험 1교시 문제지 원본.pdf
%% conversion: kordoc 4.24.51 -> tools/gen1.py (auto); review=proofread-sonnet; flags=none
\begin{question}{37}
질권에 관한 설명으로 옳지 않은 것은? (다툼이 있으면 판례에 따름)
\begin{choices}
  \item 질물을 질권설정자가 계속 점유하는 방식으로 질권을 설정할 수 없다.
  \item 채무자가 질권자의 책임전질을 승낙한 경우 채무자는 전질권자의 동의 없이 질권자에게 채무를 변제하여도 이로써 전질권자에게 대항하지 못한다.
  \item 질물의 소유자가 질물을 다른 사람에게 매도한 때에는 질권자는 특별한 사정이 없는 한 그 매매대금에 대하여 질권을 행사할 수 없다.
  \item 동산질권의 선의취득이 인정되려면 취득자는 자신의 선의ㆍ무과실을 증명해야 한다.
  \item 근질권이 설정된 금전채권에 대해 제3자의 압류로 강제집행절차가 개시된 경우 근질권의 피담보채권은 강제집행이 개시된 때에 확정된다.
\end{choices}
\end{question}
\answer{5}
